\documentclass[14pt, a4paper]{extarticle}

% Підтримка української мови
\usepackage[T2A]{fontenc}
\usepackage[utf8]{inputenc}
\usepackage[ukrainian]{babel}

% Математичні пакети (ОБОВ'ЯЗКОВО для роботи системи рівнянь з intro.tex)
\usepackage{amsmath, amssymb, amsfonts}

% Стандартні відступи для курсових (ДСТУ)
\usepackage[left=3cm, right=1.5cm, top=2cm, bottom=2cm]{geometry}

% Відступ першого рядка абзацу (1.25 см)
\usepackage{indentfirst}
\setlength{\parindent}{1.25cm}

\begin{document}
	
	\textbf{Актуальність теми.}
	Наявне наукове протиріччя полягає в гострому розриві між потребою клінічної медицини в індивідуалізації протоколів імунотерапії та неможливістю проведення in vivo експериментів для перевірки всіх комбінацій дозувань. Це зумовлює критичну необхідність розробки методів комп'ютерного моделювання.
	
	Аналіз останніх досліджень показує, що протягом останніх років над математичним описом онкоімунологічних процесів активно працювали H.~Enderling, D.~Kirschner, A.~d'Onofrio та вітчизняні фахівці, які заклали основи популяційної динаміки на базі звичайних диференціальних рівнянь.
	
	Проте невирішеною частиною проблеми залишається адаптація цих диференціальних моделей для швидкого сценарного аналізу та візуалізації результатів у сучасних програмних середовищах. Саме розробка такого програмно-математичного інструментарію, зокрема засобами мови програмування Python, для симуляції різних режимів терапії становить ключовий акцент і визначає актуальність даного дослідження.
	
	\textbf{Мета дослідження:} розробити математичну модель та її програмну реалізацію для прогнозування впливу імунотерапії на динаміку взаємодії пухлинних та імунних клітин шляхом проведення сценарного аналізу.
	
	\textbf{Об'єкт дослідження:} процес динамічної взаємодії популяцій пухлинних та імунних клітин в організмі людини в умовах зовнішнього терапевтичного впливу.
	
	\textbf{Предмет дослідження:} математичні моделі на базі нелінійних систем звичайних диференціальних рівнянь та програмні алгоритми, що кількісно описують вплив імунотерапії на динаміку взаємодії зазначених клітинних популяцій.
	
	% Підключення файлу з математичною моделлю
	\newpage
	\section{Постановка задачі}

\subsection{Біологічне обґрунтування моделі}
Проблема лікування онкологічних захворювань залишається однією з найгостріших у сучасній медицині. Традиційні методи, такі як хіміотерапія та променева терапія, мають значний токсичний вплив на здорові тканини організму. Останніми роками дедалі більшої уваги набуває імунотерапія — підхід, що базується на активації та стимуляції власної імунної системи пацієнта для розпізнавання та знищення злоякісних клітин.

Для розробки ефективних протоколів лікування необхідно розуміти складну динаміку взаємодії між пухлиною та імунітетом. Математичне моделювання дозволяє не лише описати ці процеси кількісно, але й прогнозувати результати різних терапевтичних стратегій. У даній роботі розглядається базова математична модель, що описує еволюцію двох взаємодіючих клітинних популяцій.

\subsection{Математична формалізація}
Нехай стан біологічної системи в довільний момент часу $t \ge 0$ описується двома фазовими змінними:
\begin{itemize}
	\item $T(t)$ — кількість (або концентрація) пухлинних клітин;
	\item $I(t)$ — кількість (або концентрація) активних ефекторних клітин імунної системи (наприклад, цитотоксичних Т-лімфоцитів).
\end{itemize}

Зміна цих популяцій у часі моделюється за допомогою системи нелінійних звичайних диференціальних рівнянь. Припускається, що за відсутності імунної відповіді пухлина зростає за логістичним законом. Взаємодія між імунними та пухлинними клітинами призводить до загибелі останніх. Водночас наявність пухлинних антигенів стимулює проліферацію імунних клітин, швидкість якої обмежується ефектом насичення (моделюється функцією типу Міхаеліса-Ментен).

Динаміка системи з урахуванням зовнішнього імунотерапевтичного впливу має вигляд:
\begin{equation}
	\begin{cases}
		\frac{dT}{dt} = r T \left( 1 - \frac{T}{K} \right) - \alpha T I \\[8pt]
		\frac{dI}{dt} = s + \rho \frac{T I}{\gamma + T} - \delta I + u(t)
	\end{cases}
\end{equation}

де параметри моделі є додатними константами та мають наступний зміст:
\begin{itemize}
	\item $r$ — питома швидкість природного розмноження пухлинних клітин;
	\item $K$ — максимальна ємність середовища, що обмежує розмір пухлини;
	\item $\alpha$ — інтенсивність знищення пухлинних клітин імунними клітинами при їхньому контакті;
	\item $s$ — постійна швидкість надходження нових імунних клітин в область пухлини (природний базовий імунітет);
	\item $\rho$ — максимальна швидкість стимуляції імунної системи під дією пухлинних антигенів;
	\item $\gamma$ — константа напівнасичення (порогове значення кількості пухлинних клітин, при якому швидкість стимуляції імунітету досягає половини від максимальної);
	\item $\delta$ — коефіцієнт природної загибелі (апоптозу) клітин імунної системи;
	\item $u(t)$ — функція керування, що невід'ємна при $t \ge 0$ і описує інтенсивність зовнішнього введення імунотерапевтичних препаратів.
\end{itemize}

\subsection{Формулювання задачі Коші та цілі дослідження}
Для повноти математичної постановки необхідно задати початкові умови, що відображають стан організму пацієнта на момент початку спостереження або лікування ($t = 0$):
\begin{equation}
	T(0) = T_0 > 0, \quad I(0) = I_0 > 0
\end{equation}

Головною задачею даної курсової роботи є:
\begin{enumerate}
	\item Знаходження координат стаціонарних точок системи за умови відсутності лікування ($u(t) = 0$).
	\item Аналітичне дослідження знайдених стаціонарних точок на стійкість за першим наближенням (аналіз власних значень матриці Якобі).
	\item Побудова фазових портретів системи.
	\item Чисельне моделювання динаміки взаємодії клітин при різних режимах імунотерапії (різних функціях $u(t)$) за допомогою програмних засобів мови Python для визначення оптимальних стратегій лікування.
\end{enumerate}
	
\end{document}